% Suggested prose only; not installed in the manuscript.
\section{Related work}

\paragraph{Constructive Gaussian program theory.}
Tensor Programs provides the closest mathematical foundation.
Tensor Programs III gives a constructive Gaussian source-response rule for
named matrix and transpose reuse, including singular query covariances and
causal scalar feedback \cite{yang2021tp3}. Non-Gaussian Tensor Programs supplies
the smooth raw-entry moment estimates relevant to expectation convergence
\cite{golikov2022nongaus}. Tensor Programs IVb gives a broader primitive
language, explicit backpropagation constructions and a master theorem with
finite-$L^p$ convergence under polynomial smoothness and uniform initial-input
moments \cite{yang2023tp4b}. Its fixed-input Gaussian branch covers MFP's
compiled primitive programs. Tensor Programs IV and IVb also construct
feature-learning training programs; after rescaling parameter blocks, the
maximal-update regime agrees with the normalization of our MLP examples
\cite{yang2021tp4,yang2023tp4b}.

These are constructive antecedents, with covariance and response calculations
already determined by the program. The further task addressed by MFP is to give
one explicit contract for neural derivative queries and carry it through to
symbolic output. This requires specifying the permitted matrix directions and
contractions, proving exact closure under their differentiation and updates,
retaining the distinction between moving and frozen directions, and connecting
the resulting programs to the Gaussian evaluator. The initialization theorem
then proves a more specific output form in activation-derivative moments.
These finite closure and reduction arguments have their own correctness
obligations, although they use established chain, adjoint and Stein identities.
The compiler implements this complete transformation. Its contribution should
be assessed as a calculus with a precise mathematical and executable domain,
with full credit to the probabilistic machinery on which it builds.

Recent quantitative work strengthens the conclusions on a different domain.
Agazzi, Mosig Garc\'{\i}a and Trevisan prove inverse-square-root-width
Wasserstein estimates for forward tensor programs with globally Lipschitz
nonlinearities, including a kernel-feedback extension under a joint
nondegeneracy assumption \cite{agazzi2026quantitative}. These are stronger
finite-width guarantees than MFP supplies on the common scope. Their language
does not include shared transpose calls, and its Lipschitz hypotheses do not
remain valid under general directional differentiation: even for
$\phi=\tanh$, the coordinate map $(z,v)\mapsto\phi'(z)v$ is not globally
Lipschitz. Thus a forward quantitative limit cannot automatically be applied
to every differentiated program in the present calculus.

\paragraph{Derivative hierarchies and diagrammatic calculation.}
The neural tangent hierarchy organizes repeated gradient contractions and
studies quantitative approximation of training in the tangent-kernel regime
\cite{huang2020nth}. Label-Aware NTK gives a recursive symbolic program for
that hierarchy and works out shallow-network initialization calculations to
motivate label-dependent kernels \cite{chen2020lantk}. These are direct
predecessors for derivative recurrences and symbolic label dependence. A
hierarchy identity specifies the next finite derivative; obtaining its
initialization expectation still requires a probabilistic calculation. MFP
supplies that calculation for its declared smooth program class, while its
finite initial jets address a different temporal claim from positive-time
hierarchy truncation.

Feynman-diagram methods organize contracted neural derivative tensors and
finite-width corrections \cite{dyer2019feynman}. The polynomial-activation
scaling theorem of Aitken and Gur-Ari supports exact Wick calculations for
arbitrary fixed derivative contractions in its model
\cite{aitken2020nonlinear}. Such calculations can recover corresponding
polynomial MFP coefficients once parameter metrics and width factors are
matched. They also retain subleading finite-width information absent from
MFP's leading-limit graph. MFP's output instead retains smooth activation
functions symbolically through their Gaussian derivative moments and is
specified for a compositional derivative/update language.

Guillen, Misof and Gerken provide an especially close computational precedent:
a diagrammatic calculus for preactivation, NTK, dNTK and ddNTK statistics,
with implemented Gaussian recursions in NTK Unlimited \cite{guillen2026finite}.
Their all-order statement concerns powers of $1/n$ in the statistics of the
designated tensors. The inspected implementation takes supplied tensor
recurrences, simplifies them symbolically and compiles them to numerical
Gaussian integration. MFP accepts the neural program and physical derivative
request itself, and derives the corresponding Gaussian graph. This distinction
concerns the mathematical input and the work automated, rather than whether
both systems contain symbolic expectation nodes. Their finite-width expansion
also develops corrections around the tangent-kernel scaling; MFP's featured
jets retain the mean-field metric in which hidden features can move at leading
order. Neither a change of names nor a finite truncation in $1/n$ identifies
those training regimes.

\paragraph{Nonlinear dynamical descriptions.}
Dynamical mean-field theory provides a different route to calculating feature
learning. Bordelon and Pehlevan formulate coupled Gaussian driving processes,
response functions and time-dependent kernels for nonlinear training, and
supply a numerical self-consistency procedure \cite{bordelon2023dmft}.
Their framework retains the response to matrix reuse and must be distinguished
from a gradient-independence approximation. Its discrete-time connection to
Tensor Programs gives a rigorous route for fixed admitted training histories;
the continuum path-integral derivation and its discrete-to-continuous limit
have separate, explicitly unresolved proof obligations. Fixed initial jets
can already be reached by differentiating a finite update program before
applying the program theorem. MFP makes that procedure uniform and executable,
with the stated moment guarantees and initialization reduction.
The difference relevant to the dynamical solver is the
computational question: solving the coupled history-dependent dynamics versus
compiling a specified finite derivative query into a causal expectation graph.
MFP's initial-jet calculation determines its coefficients without solving for
an unknown positive-time kernel trajectory. It does not replace a dynamical
solution or assert that a finite list of initial coefficients reconstructs one.

\paragraph{Software for kernels, derivatives and Gaussian algebra.}
Neural Tangents automates supported architecture-to-NNGP/NTK calculations,
provides numerical kernel evaluation and prediction tools, and supports
arbitrary-order finite-network Taylor expansions
\cite{novak2020neuraltangents}. Tensorgrad supplies symbolic tensor
differentiation, Gaussian expectations with symbolic means and covariances,
and Wick/Stein simplification \cite{ahle2026tensorgrad}. These are substantial
capabilities: higher-order differentiation and Gaussian symbolic integration
are already available, including compositions that reproduce elementary
moving-derivative examples.

MFP's distinctive executable task is the elimination of growing-width neural
matrix calculations through the response-corrected program limit. Its general
moving-jet operation, represented updates, formal-source bookkeeping and
initialization reduction act on the same typed input language. Among the
interfaces inspected, we found no earlier package exposing that complete
neural-program-to-derivative-limit transformation. This bounded comparison
supports the stated contribution without requiring priority for its individual
operations. Existing software also supplies capabilities beyond MFP, including
numerical quadrature, wider architecture support and broader tensor algebra.
The relevant test is which complete calculation each system accepts and returns.
